\section{Conclusão e Perspectivas Futuras}

A intersecção e sinergia entre a implantodontia cirúrgica e a inteligência computacional consubstanciada no Gêmeo Digital demarca, de modo categórico, uma metamorfose no planejamento, execução tangível e monitoramento remoto de reabilitações craniofaciais em osso vivo. Esta transmutação em direção a ecossistemas interconectados — balizados pela fusão entre macro e microgeometrias, \textit{machine learning} acoplado à mecanobiologia, manufatura aditiva submilimétrica e arquiteturas IoT de sensoriamento contínuo ininterrupto — assegura graus exatos de precisão que subvertem as limitações da dependência heurística secular do operador clínico~\cite{khoshfekr2025digital, duggal2026exploring}.

Como atestado detalhadamente nas simulações fenotípicas exploradas (da exodontia preditiva orientada por vetores de escoamento tensão-deformação~\cite{xing2024clinical} às complexidades bioengenheiradas de texturização de superfície não colinear~\cite{alshadidi2024surface}), a robustez destas metodologias atua no substrato de prever em silício o que se consumará em material genético celular biológico. Contudo, para a plena fruição de uma odontologia de \textit{Big Data} irrestrita, obstáculos ontológicos e socioeconômicos basilares devem ser vencidos em caráter mandatório~\cite{du2020marginalized}. Ensaios clínicos independentes e não patrocinados que confrontem desfechos longitudinais com padrões clássicos, alinhados com o redesenho curricular universitário profundo, são prerrogativas cardeais.

A projeção imanente para as décadas subsequentes da medicina orofacial pauta o refinamento das redes semânticas e do aprendizado federado (\textit{Federated Learning}), possibilitando que o Gêmeo Digital deixe de atuar meramente como emulador reativo cirúrgico para assumir a governança como um autêntico sistema de suporte preditivo omnisciente contínuo e cibernético~\cite{techsoc2025precision}. Com a adoção global de interfaces cognitivas estandardizadas e protocolos descentralizados e abertos (\textit{Open-Source}), como o promissor ecossistema em nuvem delineado no movimento \textit{OpenTwins}~\cite{robles2023opentwins}, o Gêmeo Digital se consolidará incontestavelmente como a fundação irrefutável do conhecimento mecatrônico, biomecânico e cirúrgico no espectro maxilofacial.
